\documentclass[11pt,onecolumn]{IEEEtran}

% =========================================================
% Packages
% =========================================================
\usepackage{amsmath}
\usepackage{amssymb}
\usepackage{graphicx}
\usepackage{enumitem}
\usepackage{xcolor}
\usepackage{titlesec}
\usepackage{hyperref}
\usepackage{newtxtext}
\usepackage{newtxmath}

% =========================================================
% Formatting
% =========================================================
\setlength{\parindent}{0pt}
\setlength{\parskip}{3pt}

\setlist[itemize]{
    leftmargin=1.5em,
    itemsep=4pt,
    topsep=3pt
}

% Plain, left-aligned run-in headings (memo style, matching the
% example proposal) instead of IEEEtran's default centered small caps.
\titleformat{\section}{\normalfont\large\bfseries}{}{0pt}{}
\titlespacing*{\section}{0pt}{14pt}{6pt}

% Subtitle color matching the example proposal's link-style byline.
\definecolor{subtitleblue}{HTML}{1155CC}

% =========================================================
% Document
% =========================================================
\begin{document}

% =========================================================
% Title
% =========================================================

\begin{flushleft}
    {\huge \textbf{Learning Planar Object Pushing with Reinforcement Learning}}\\[6pt]
    {\large \color{subtitleblue} ME5418 Final Project Proposal}\\[10pt]

    \textbf{Taowen Liu}, A0352191J
    \qquad
    \textbf{Xuanwei Liu}, A0357433B
    \qquad
    \textbf{Chengtao Zheng}, A0351467Y
\end{flushleft}

\vspace{0.2cm}


% =========================================================
% Why
% =========================================================

\section*{Why:}

\begin{itemize}

    \item
    Pushing is a fundamental robotic manipulation skill that can be useful when
    directly grasping an object is difficult or unnecessary. For example, in
    industrial manipulation, a robot may first push a component to a desired
    position and orientation before grasping, assembling, or further manipulating
    it.

    \item
    Despite its simple appearance, pushing is a contact-rich manipulation problem.
    The resulting object motion depends on the contact location, pushing direction,
    object geometry, and physical properties such as friction. It can therefore be
    difficult to design a single hand-crafted controller that works reliably under
    different initial configurations.

    \item
    In this project, we investigate whether reinforcement learning can learn a
    closed-loop pushing policy that moves an object from a randomized initial pose
    to a desired target pose. We will primarily consider a T-shaped object similar
    to the Push-T task, while the same framework can also be evaluated on simpler
    object geometries.

\end{itemize}


% =========================================================
% Conventional Algorithms
% =========================================================

\section*{Conventional Algorithms:}

\begin{itemize}

    \item
    Conventional robotic pushing approaches can use geometric planning,
    model-based control, trajectory optimization, or manually designed pushing
    strategies. Given an estimate of object pose and contact dynamics, these
    approaches can determine a sequence of pushing actions to move the object
    toward a target configuration.

    \item
    However, accurate contact dynamics can be difficult to model, especially when
    object geometry and friction vary. Model-based approaches may therefore require
    accurate physical parameters or repeated replanning.

    \item
    Instead, we would like to investigate whether a reinforcement learning agent
    can directly learn a feedback policy through interaction with the simulated
    environment, and adapt its pushing actions according to the current object
    and target poses.

\end{itemize}


% =========================================================
% Problem Statement
% =========================================================

\section*{Problem Statement:}

\begin{itemize}

    \item
    We consider a tabletop planar pushing environment containing a robot
    end-effector (or pusher), a movable rigid object, and a target pose. The main
    experiment will use a T-shaped object, although simpler shapes may also be
    considered. The workspace is a bounded square tabletop (e.g., roughly
    $0.4\,\mathrm{m} \times 0.4\,\mathrm{m}$), simulated with a 2D rigid-body
    physics engine (e.g., PyMunk or a similar lightweight simulator, or MuJoCo if
    a 3D pusher/arm is used).

    \item
    At the beginning of each episode, the object's initial position and orientation
    are randomly generated within the workspace. The target position and
    orientation may also be randomized.

    \item
    The robot must push the object from its initial configuration to the target
    configuration within a limited number of control steps (e.g., 200--300 steps
    per episode). The task is completed successfully when both the position and
    orientation errors are below predefined thresholds, e.g., a position error
    below $2\,\mathrm{cm}$ and an orientation error below $10^{\circ}$.

    \item
    If time permits, we will additionally randomize physical properties such as
    the object--table friction coefficient to test whether the learned policy can
    remain robust under different dynamics.

\end{itemize}


% =========================================================
% RL Cast
% =========================================================

\section*{RL Cast:}

\begin{itemize}

    \item
    \textbf{State space:}

    The observation will contain low-dimensional robot and object states,
    including:

    \begin{itemize}
        \item current pusher/end-effector position;
        \item object position and orientation;
        \item target position and orientation;
        \item relative object-to-target position and orientation;
        \item optionally, object and robot velocities.
    \end{itemize}

    \item
    \textbf{Action space:}

    We will use a continuous action space controlling the planar motion of the
    pusher/end-effector. A simple action representation is

    \[
        a_t = [\Delta x, \Delta y],
    \]

    where $\Delta x$ and $\Delta y$ specify the desired displacement of the
    pusher during each control step.

    If a robotic arm is used, these commands can instead be converted into
    end-effector Cartesian control commands.

    \item
    \textbf{Reward structure:}

    The reward will mainly encourage progress toward the desired object pose.
    A possible reward is

    \[
        r_t =
        -w_p d_{\mathrm{pos}}
        -w_r d_{\mathrm{rot}}
        - c_{\mathrm{step}}
        + r_{\mathrm{success}},
    \]

    where $d_{\mathrm{pos}}$ is the distance between the current and target object
    positions, $d_{\mathrm{rot}}$ is the orientation error, $c_{\mathrm{step}}$ is
    a small constant per-step cost, and $r_{\mathrm{success}}$ is a large positive
    reward when the object reaches the target pose. As a starting point we plan to
    use $w_p \approx 1.0$, $w_r \approx 0.5$, $c_{\mathrm{step}} \approx 0.01$, and
    $r_{\mathrm{success}} = +10$; these weights will be tuned empirically.

    An additional penalty (e.g., $-1$) may be applied if the pusher exits the
    workspace or loses contact with the object for an extended period, to
    discourage degenerate behaviors.

    \item
    \textbf{Neural Network Structure:}

    Since the observation space is low-dimensional, we plan to use a simple
    multilayer perceptron (MLP) for the policy and value networks, e.g., two
    hidden layers of 256 units with ReLU activations, following common defaults
    for continuous-control PPO implementations (such as Stable-Baselines3). The
    actor and critic may share the same trunk with separate output heads, or use
    fully separate networks if training is unstable.

\end{itemize}


% =========================================================
% RL Algorithms
% =========================================================

\section*{RL Algorithms and Envisioned Results:}

\begin{itemize}

    \item
    We plan to use Proximal Policy Optimization (PPO) as the primary reinforcement
    learning algorithm, since it is widely used for continuous robotic control and
    is relatively stable and straightforward to train. Depending on the available
    time, we may additionally compare PPO with another continuous-control
    algorithm such as SAC.

    \item
    We will evaluate the learned policy mainly in terms of task success rate,
    final position error, final orientation error, and the number of control steps
    required to reach the target, averaged over a held-out set of randomized
    episodes (e.g., 100 test episodes not seen during training).

    \item
    The learned policy will be evaluated under different object initial poses and
    target poses. If time permits, we will further test different friction
    coefficients or object geometries to study the generalization and robustness
    of the policy.

    \item
    As a simple baseline, we will compare the RL policy with a heuristic pushing
    controller that first moves the pusher to a contact point behind the object
    (relative to the target direction) and then pushes it approximately in a
    straight line toward the target, re-planning the contact point whenever the
    object-to-target direction changes significantly. We expect the learned
    policy to perform better on cases that require multiple contacts or
    simultaneous adjustment of object position and orientation, where the
    straight-line heuristic is expected to fail or take a large number of steps.

\end{itemize}


\end{document}
